\section{Related Work}

\label{sec:related}
% ============================================================

\paragraph{Bridge Governance.}
Wormhole~\cite{wormhole} uses a guardian council of 19 entities for both security and governance. LayerZero~\cite{layerzero} separates message validation (Oracle + Relayer) from application logic but does not provide per-chain fee governance. Axelar~\cite{axelar} uses delegated Proof-of-Stake for governance, concentrating control in token holders of the Axelar chain rather than connected chains. Teleport's sovereign model is closest to IBC~\cite{ibc}'s design, where each chain controls its own IBC module, but extends to non-Cosmos chains.

\paragraph{Islamic Finance and DeFi.}
Mrhool~\cite{mrhool2022} surveys Shariah compliance in DeFi and identifies the absence of protocol-level enforcement as the primary barrier. ShariaChain~\cite{shariachain2021} proposes a Shariah-compliant blockchain but does not address cross-chain compliance. HAQQ Network~\cite{haqq2023} implements Islamic finance principles at the L1 level but does not provide programmable yield classification. Teleport's ShariaFilter is the first mechanism to enforce Shariah compliance at the bridge protocol layer, enabling any chain to opt into compliance without modifying its own consensus or execution environment.

\paragraph{Security Token Standards.}
ERC-3643 (T-REX)~\cite{erc3643} defines an on-chain identity and compliance framework for security tokens. Polymath's ST-20~\cite{polymath2018} provides a similar framework. Neither standard addresses cross-chain compliance preservation. Teleport's securities bridge mode extends ERC-3643 across chain boundaries through MPC-attested registry synchronization. The EU MiCA regulation~\cite{mica2023} establishes requirements for crypto-asset service providers that Teleport's settlement registry is designed to satisfy.

\paragraph{Fee Mechanism Design.}
Roughgarden~\cite{roughgarden2021} analyzes transaction fee mechanisms for blockchains. Our fee competition model extends this to the multi-instance bridge setting, where sovereign instances compete on fee and liquidity depth rather than block space. The hard cap at 100~bps ensures that fee competition cannot produce welfare-destroying outcomes.

% ============================================================
